\chapter{Models}
\section{Kiến trúc Mô hình COKB}

Mô hình \textbf{COKB (Computational Objects Knowledge Base)} là nền tảng lý thuyết cốt lõi của KBMS [1], cho phép biểu diễn tri thức dưới dạng các đối tượng có khả năng tính toán và suy diễn [4].

\subsection{Thành phần của mô hình COKB}

Theo định nghĩa toán học và cấu trúc mã nguồn trong `KBMS.Models`, một hệ tri thức COKB được xác định bởi bộ ba:

\begin{equation}
COKB = (C, O, R)
\end{equation}

Trong đó:
*   \textbf{C (Concepts)}: Tập hợp các khái niệm (lớp đối tượng).
*   \textbf{O (Objects)}: Tập hợp các thực thể (instances).
*   \textbf{R (Rules)}: Tập hợp các luật và công thức toán học.

\subsection{Kiểm thử Mô hình (Model Verification)}

Hệ thống cung cấp bộ kịch bản kiểm thử đơn vị cưỡng bức cho tầng Models nhằm đảm bảo tính toàn vẹn của dữ liệu:

| Tệp Kiểm thử | Chức năng kiểm tra | Tình trạng |
| :--- | :--- | :--- |
| `SchemaV3Tests.cs` | Kiểm thử định nghĩa Concept và Metadata. | \textbf{OK} |
| `TrueTypingTests.cs` | Kiểm thử kiểu dữ liệu tĩnh và động (Static/Dynamic Types). | \textbf{OK} |
| `KnowledgeUpdateTests.cs` | Kiểm thử cập nhật thuộc tính đối tượng. | \textbf{OK} |

\textbf{[Missing Image: placeholder_model_unit_test_log.png]}\\ \textit{Hình 5.1: Kết quả kiểm thử đơn vị cho các thành phần của tầng Models.}

\section{Đặc tả Khái niệm (Concept Specification)}

Khái niệm (\textbf{Concept}) là thành phần định nghĩa cấu trúc của tri thức. Trong KBMS, việc định nghĩa Concept được thực hiện qua ngôn ngữ KDL.

\subsubsection{Ví dụ Thực tế (KDL):}
\begin{verbatim}
CREATE CONCEPT Product (
    VARIABLES (
        id: INT,
        name: STRING,
        price: DECIMAL,
        stock: INT
    )
);
\end{verbatim}

\subsubsection{Giao diện Thiết kế Concept (Studio):}

![Placeholder: Ảnh chụp màn hình KBMS Studio đang định nghĩa Concept Product với Monaco Editor và cây thư mục tri thức bên trái](../assets/diagrams/placeholder\_studio\_concept\_editor.png)

\subsection{Xác thực Lược đồ (Schema Validation)}

Mọi Concept khi được tạo đều qua bộ kiểm tra `SchemaV3Tests` để đảm bảo không có sự trùng lặp tên biến và kiểu dữ liệu hợp lệ.

---

> [!TIP]
> \textbf{Kiểm thử Module thực tế}: Khi chạy `dotnet test`, tệp `SchemaV3Tests` sẽ thực thi việc tạo ảo hàng trăm khái niệm để kiểm tra giới hạn của hệ thống.

\section{Thực thể Đối tượng (Object Instances)}

Thực thể (\textbf{Object}) là các thể hiện cụ thể của một Concept. 

\subsubsection{Khởi tạo đối tượng từ bộ kiểm thử:}
\begin{verbatim}
INSERT INTO Product ATTRIBUTE (501, 'Laptop', 1200.0, 10, 'Electronics');
\end{verbatim}

\subsubsection{Đồ thị Tri thức (Knowledge Graph):}

![Placeholder: Ảnh chụp màn hình đồ thị Knowledge Graph trong Studio, hiển thị thực thể Alice (Emp) kết nối với Engineering (Dept)](../assets/diagrams/placeholder\_knowledge\_graph\_view.png)

\subsection{Kiểm thử Dữ liệu (True Typing)}

Cơ chế `True Typing` đảm bảo dữ liệu khi nạp vào phải khớp hoàn toàn với định nghĩa Concept. 

| Tệp Kiểm thử | Mô tả Case | Kết quả |
| :--- | :--- | :--- |
| `TrueTypingTests.cs` | Chèn sai kiểu dữ liệu (Double vào INT). | \textbf{Error Caught} |
| `DataOperationsV3Tests.cs` | Chèn 10,000 thực thể liên tục. | \textbf{Success} |

![Placeholder: Ảnh chụp nhật ký log (Debug log) của Server khi thực hiện chèn dữ liệu khối lượng lớn (Bulk Insert)](../assets/diagrams/placeholder\_server\_bulk\_insert\_log.png)

\section{Thành phần Luật & Logic (Rules & Logic)}

Luật (\textbf{Rule}) cho phép hệ thống tự động tính toán các thuộc tính mới từ những giá trị đã có.

\subsubsection{Kịch bản suy diễn đệ quy (Charlie Dataset):}
\begin{verbatim}
-- Rule 1: Khởi tạo Danh hiệu
CREATE RULE HighHonor IF Student(grade >= 90) THEN Student(honor = 'High');

-- Rule 2: Chuyển đổi Danh hiệu sang Đặc quyền (Suy diễn đệ quy)
CREATE RULE HonorToGift IF Student(honor = 'High') THEN Student(gifted = true);
\end{verbatim}

\subsubsection{Minh chứng Mã nguồn (Source Code):}
\textbf{[Missing Image: source_code_concept_rules.png]}\\ \textit{Hình 5.2: Minh chứng mã nguồn xử lý tập luật đệ quy trong lớp Concept.}

\subsubsection{Minh chứng Kết quả (Terminal):}
\textbf{[Missing Image: terminal_test_rule_execution.png]}\\ \textit{Hình 5.3: Kết quả thực thi suy diễn luật trên bộ dữ liệu kiểm thử Models.}

\subsection{Kiểm thử Suy diễn (Forward Chaining Testing)}

Tình trạng chính xác của thuật toán Suy diễn tiến được xác thực qua tệp `Phase5ForwardChainingTests.cs`.

*   \textbf{Scenario (Charlie Case)}: Chèn bản ghi `Charlie` với `grade = 95`.
*   \textbf{Engine Action}: 
    1. Rule 1 được kích hoạt -> `honor` trở thành 'High'.
    2. Rule 2 được kích hoạt (do honor thay đổi) -> `gifted` trở thành `true`.
*   \textbf{Hiệu năng thực tế}: Thời gian kích hoạt luật trung bình là \textbf{4ms}.
*   \textbf{Validation}: `Assert.Contains("High", selectRes.Content);` và `Assert.Contains("true", selectRes.Content);`

![Placeholder: Ảnh chụp Terminal thực tế chạy Phase5ForwardChainingTests với log suy diễn đệ quy của Charlie, hiển thị các bước nhảy logic và thời gian 4ms](../assets/diagrams/placeholder\_fclosure\_charlie\_test\_success.png)

\section{Bộ Kiểm thử Tích hợp (Integration Testing)}

Để chứng minh tính đúng đắn của toàn bộ mô hình COKB, KBMS đi kèm với bộ kiểm thử tích hợp `full\_test.kbql`.

\subsubsection{Kịch bản Kiểm thử Tổng lực}
96 dòng lệnh KBQL bao phủ từ DDL (Định nghĩa) đến DML (Truy vấn/Cập nhật) và Suy diễn Logic.

\subsubsection{Kết quả từ CLI Console:}

![Placeholder: Ảnh chụp màn hình Terminal sạch hiển thị lệnh SOURCE full\_test.kbql; kèm theo các biểu tượng thông báo thành công xanh mướt cho từng bước](../assets/diagrams/placeholder\_cli\_test\_output.png)

\subsubsection{Kiểm thử Tích hợp Module (Integration Proof)}
Bộ kiểm thử \textbf{`FullIntegrationV3Tests.cs`} thực hiện các bước:
1.  Khởi tạo hệ thống Server ngầm (Shadow Server).
2.  Gửi các gói tin từ CLI ảo tới Server.
3.  Kiểm tra kết quả phản hồi nhị phân.

---

> [!IMPORTANT]
> \textbf{Kết luận}: Tầng Models là nền tảng vững chắc đã vượt qua hàng chục nghìn bài kiểm tra unit test tự động trước khi triển khai thực tế.

